\subsection{The dimension order and signed central traces}\label{sec:dimension-order}
The positive-dimension induction system has two extensions that are useful
when an exhaustive sphericalization search leaves a non-positive
spherical character. The first sharpens an arithmetic condition without
adding any central variables. The second retains the algebraic part of
the same trace system, but discards every inequality that depends on
positivity of categorical dimensions.

\subsubsection{The actual order of simple dimensions}
Let $d=(1,d_1,\ldots,d_{r-1})$ be a dimension character under
consideration, and let $K=\QQ(d_i)$ and $\delta=\sum_i d_i d_{i^*}$.
The finitely generated group
\[
 \mathcal O_d=\sum_i\ZZ d_i\subset K
\]
is an order: it contains $1$, consists of algebraic integers, and is
closed under multiplication because the $d_i$ satisfy the multiplication
equations with integer coefficients. It need not be the maximal order
of $K$.

\begin{proposition}[Dimension-order obstruction]\label{prop:dimension-order}
If a fusion ring $R$ has a spherical categorification with simple
dimensions $d_i$, then for every formal codegree $f$ of $R$,
\begin{equation}\label{eq:dimension-order}
 f\in K,\qquad \delta/f\in\mathcal O_d.
\end{equation}
The same conclusion holds for the dimensions of selected lifts in the
canonical sphericalization of an arbitrary categorification of $R$.
Thus the order-membership test may be applied to each dimension-character
branch of a complete signed-sphericalization search.
\end{proposition}
\begin{proof}
In the spherical case the central simple $A_E$ attached to an irreducible
representation of $R$ with codegree $f$ has dimension $\delta/f$, by
\cite[Theorem~2.13]{Ostrik3}. Its forgetful image is an object of the
original category. Its dimension is therefore a nonnegative integral
linear combination of the $d_i$, in particular an element of
$\mathcal O_d$. Since $\delta$ is nonzero, this also puts $f$ in $K$.

For an arbitrary categorification use its sphericalization, of dimension
$2\delta$. An irreducible representation of $R$ lifts through the
forgetful quotient to one with formal codegree $2f$; both selected and
opposite basis lifts contribute to that codegree. The same central
induction theorem gives dimension $(2\delta)/(2f)=\delta/f$.
All simple dimensions of the sphericalization are among $\{d_i,-d_i\}$,
so forgetting this central object gives an integral linear combination
of the selected dimensions. This proves the assertion without assuming
that the selected lifts themselves form a fusion subcategory.
\end{proof}

The test is finite and exact. Express the $d_i$ in a rational power basis
of $K$, clear one common denominator, and put their coordinate lattice
in Hermite normal form. Membership of $\delta/f$ is then an integer
linear-algebra question. A negative certificate can equivalently give a
rational linear functional $\ell:K\to\QQ$ satisfying
\[
 \ell(d_i)\in\ZZ\quad\hbox{for all }i,
 \qquad \ell(\delta/f)\notin\ZZ.
\]
The checker verifies the defining polynomial of $f$, reconstructs the
lattice and the functional, and checks both displayed conditions. A
failure on one character component is not a noncategorification proof
unless all possible sphericalization and character branches are closed.
In particular a failed character of a positively realized ring cannot
be promoted to an exclusion of that ring.

For illustration, the complete certificate for $\RR{4}{375}$ has four
oriented signed tables, all gauges of the original tensor. Its nonzero
spherical components have order $\ZZ[\sqrt5]$ and dimensions
\[
 \delta=70+30\sqrt5\quad\hbox{or}\quad\delta=30-10\sqrt5.
\]
The codegree $4$ gives quotients outside that order in both cases.
The supplied functional certificates, not this identifier or this
displayed calculation, are the inputs to the automated verification.

\subsubsection{The algebraic trace system without positive dimensions}
Suppose a spherical character takes values in a real field $K$ of degree
at most two, has nonzero simple dimensions, and satisfies
$d_i=d_{i^*}$. Write $\delta=\sum_i d_i^2$. The restriction vectors of
central simples remain nonnegative integral vectors, even when some
categorical dimensions are negative. Consequently the Gram identity
\eqref{eq:fullgram}, its finite entry bounds, and the formal-codegree
low-column equalities remain valid. They do not use an inequality
between categorical dimensions.

For a commutative input let $\chi_a$ be a character label of a low simple
$Z_a$ and put $\beta_a=\delta/f_a$. If $v_W$ is a candidate restriction
vector, define
\[
 d_W=d\cdot v_W,\qquad
 H_{aW}=\beta_a\chi_a(v_W),\qquad
 \Lambda_W=(H_{aW}/d_W)_a.
\]
Assume that the full ring character table is defined over $K$. The
character correspondence of Theorem~\ref{thm:character-induction} gives
$H_{aW}=S_{Z_a,W}$, with $S$ unnormalized. The following formulation
uses no absolute value of $d_W$.

\begin{theorem}[Signed central-induction obstruction]\label{thm:signed-traces}
A spherical categorification with the specified character supplies a
nonnegative integral solution of the following finite necessary system.
The unknowns count high central restriction vectors and their joint
relative twist-trace types. They satisfy the Gram equations, the signed
identities
\begin{equation}\label{eq:signed-indicators}
 \sum_W(v_W)_i d_W\theta_W^j=\delta\nu_j(b_i),
\end{equation}
and the cyclic-operator bounds on the right-hand indicators. Low
columns have twist one and dimension $\beta_a$. They admit at least
one complete character labelling satisfying modular reciprocity.

All coordinates of $\Lambda_W$ are algebraic integers. Their joint
Galois types have balanced multiplicities, separately for each relative
twist-trace type. An absent Galois type forbids the corresponding
restriction vector. Finally, if $P$ is the normalized trace from a
common cyclotomic field to $K$, then
\begin{equation}\label{eq:signed-low-modular}
 \sum_W P(\theta_W)H_{aW}H_{bW}=\delta S_{Z_a,Z_b}
\end{equation}
for every pair of low labels.

Exhaustive failure of the system for all bounded low columns, all
compatible character labellings, and all permitted second indicators
excludes the specified spherical character. Exact rational Farkas
certificates suffice for failure of the nonnegative rational relaxation.
\end{theorem}
\begin{proof}
Adjunction proves the Gram identity with nonnegative restriction
multiplicities. Central induction supplies the low dimensions and their
unit multiplicities exactly as in \eqref{eq:weightedlow}; neither
statement assumes pseudo-unitarity. The trace formula
\cite{NS} proves \eqref{eq:signed-indicators} with the actual, possibly
negative, dimensions. Bounds obtained as traces of finite-order cyclic
operators on ordinary complex vector spaces do not assume positivity
of categorical dimensions. Only these algebraic bounds, and not the
positive low-mass inequalities, are retained here.

The finite relative-trace list is justified by the exponent-prime
support theorem \cite{BNRW} and the same relative-trace calculation
used in \eqref{eq:orderbox}, with $\delta$ in place of $D$. The high
columns are bounded by the positive integral Gram matrix, not by the
sign of $d_W$. Their additional algebraic-integrality and $d$-number
conditions follow from the spherical modular center as in the proof of
Theorem~\ref{thm:balanced}. Columns of dimension zero are impossible.

Each $\Lambda_W$ is the restriction of the simultaneous character of
integer fusion matrices of the center to the labelled low objects.
It therefore consists of algebraic integers, and its full Galois orbit
occurs with the corresponding multiplicities. For the normalized
center representation $S/\delta,T$, the permutation of columns sends
$\theta_W$ to $\sigma^2(\theta_W)$; see
\cite[Theorem~2.9]{BNRW}. A square fixes the real field $K$ of degree
at most two, so normalized relative traces of every twist power are
unchanged. This proves the asserted balances of joint eigenvalue types.
It does not identify $d_W$ with its absolute value, nor impose the
positive dimension-orbit rule \eqref{eq:augmented-galois}.

The center has Gauss sum $\delta$. Its modular relation, together with
the twist-one property of low objects, is
$\sum_W\theta_W S_{Z_a,W}S_{W,Z_b}=\delta S_{Z_a,Z_b}$.
All coefficients other than the twists lie in $K$. Taking $P$ proves
\eqref{eq:signed-low-modular}. The actual central simples give integer
counts satisfying all the displayed equations. Expanding them in a
rational basis of $K$ and relaxing to nonnegative rational counts can
only enlarge the solution set. Farkas' lemma therefore proves the
last assertion.
\end{proof}

\paragraph{Exact scope of the implementation.}
The routine first checks the complete signed character table, its
multiplication equations, duality, nonzero dimensions and field degree.
It enumerates low columns using Gram bounds and exact dimension
equalities, without calling a positive-dimension low-mass test. It then
uses signed versions of the inherited indicator rows and adds joint
spectrum and low modular rows. In particular it does not reuse the
absolute-value dimension balances of
Proposition~\ref{prop:modular-frontier}. Only the second-indicator rows depend on the chosen signs, and each such row depends on one sign. The checker reconstructs the all-minus and all-plus endpoint systems and selects the exact normalized row for each sign choice. Independent mixed-sign reconstructions check this optimization; every separating-vector inequality is still verified. The second indicators of
self-dual objects are both allowed unless an independently checked
condition forces their sign.

For $\RR{4}{446}$ the remaining non-positive component is
\[
 d=(1,1,-1,-1,2-\sqrt6,2+\sqrt6),\qquad\delta=24.
\]
There are $290$ low matrices. Reciprocity leaves two, each with two
character labellings; all $32$ choices of second indicators are retained
for each labelling. The $128$ exact separation certificates exclude
these systems. The positive character component has the preceding
character-induction obstruction. The full oriented search has sixteen
gauge tables, so the assembled certificate excludes general
categorification, not just this displayed spherical structure.

\subsubsection{Application and structure-specific conclusions}
Component-level rejection and general nonexistence are different
conclusions. For example, rejecting one character of $\RR{4}{174}$ does
not exclude that ring. Complete proof trees for the order and signed
framework include
\[
 \RR{4}{375},\quad\RR{4}{446},\quad
 \RR{5}{1507},\quad\RR{5}{1508},\quad\RR{5}{1582},\quad
 \RR{6}{2767},\quad\RR{7}{1907}.
\]
For the last five, shorter norm and moment certificates close the
exhaustive signed alternatives, and these are retained rather than
credited to a stronger criterion unnecessarily. The finite applications
are specified by full tensors and exact proof trees.

The shorter pseudo-unitary certificates for $\RR{7}{3582}$ and
$\RR{7}{3812}$ are intermediate ingredients. Their unrestricted
existence branches are closed by the ramified signed reductions and
Casimir trace criteria above; both rings are generally excluded.

\subsection{Sphericality and Galois reduction at multiplicity three}
\label{sec:frontier-spherical-reduction}
\begin{proposition}\label{prop:frontier-galois-pu}
Every complex fusion category with based ring $\RR{3}{215}$ admits a
spherical structure and has a pseudo-unitary Galois conjugate.
Consequently general categorifiability of this ring is equivalent to
pseudo-unitary categorifiability. This reduction alone asserts neither
existence nor nonexistence; the final obstruction is proved in
Theorem~\ref{thm:frontier-noncategorifiable}.
\end{proposition}
\begin{proof}
Apply the exhaustive oriented signed-sphericalization construction of
Section~\ref{sec:signed}, using every coefficient in its full parity
interval and imposing no positive-indicator or modulo-four normalization.
There are $1480$ associative signed quotients. The $32$ unit-preserving
sign gauges partition them into $63$ orbits; one orbit consists of the
$32$ gauges of the original fusion tensor. On representatives of the
other $62$ orbits there are $271$ rational character components.
Of these, $228$ have a zero simple dimension and are impossible.
For each of the other $43$, the norm of the prospective global
dimension has $17$-adic valuation zero. The original ring has a formal
codegree with minimal polynomial $t^2-272t+1088$, whose norm has
$17$-adic valuation one. Thus the signed-sphericalization
norm-divisibility condition rejects every remaining component. No
nongauge quotient can occur for an actual category.

For completeness, a gauge quotient gives sphericality of the original
category, not merely compatibility with spherical dimensions. If
$T_{ij}^k=\eta_i\eta_j\eta_kN_{ij}^k$, replace the chosen lift of $X_i$
in the sphericalization by its tensor product with the sign object
when $\eta_i=-1$. The new signed quotient has $T=N$. The sum and
difference of the multiplicities of the two lifts of $X_k$ in a tensor
product are then both $N_{ij}^k$, so the opposite lift has multiplicity
zero. These chosen lifts therefore generate a full tensor subcategory
of the sphericalization which projects tensor-equivalently onto the
original category. Transport its spherical structure along that
equivalence.

The four rational characters of the original ring have a zero entry,
so none is a spherical dimension character. The only remaining
character orbit is
\[
 (1,1,d,d,d-1,d+1),\qquad d^2-8d-1=0.
\]
One of its two real embeddings is the Frobenius--Perron character,
$d=4+\sqrt{17}$. Galois conjugation to that embedding gives a spherical
category with positive Frobenius--Perron dimensions, hence a
pseudo-unitary category. Conversely a pseudo-unitary categorification
is a complex fusion categorification.

The complete unrestricted signed enumeration has $274655$ search nodes
and is frozen in
\filename{research/new_frontier/frontier_spherical_reduction.json}.
The companion Python script recomputes all $1480$ quotients and the
complete sign orbits, verifies the original gauge orbit, checks all
$271$ nongauge character components and their arithmetic obstructions,
and reconstructs the original dimension-character orbit using exact
arithmetic. The older normalized twelve-table enumeration is retained
separately as a regression control; it is not needed for this proof.

\end{proof}


